% ============================================================
% example-report.tex — harryopo-report.cls 完整示例
% 展示：封面、目录、章节结构、摘要、图表、参考文献
% ============================================================
\documentclass{harryopo-report}

% --- 报告信息 ---
\title{基于深度学习的深圳智慧交通流量预测系统\\研究与实现}
\subtitle{——以深圳市福田区为例}
\author{张三\quad 李四\quad 王五}
\institute{计算机学院}
\date{2026年6月21日}

% --- 基金信息 ---
\fundproject{基金项目：深圳市科技计划项目（No. JCYJ20240001）；科研基金（No. harryopo2025ZD03）}

\begin{document}

% ==== 封面 ====
\maketitle

% ==== 摘要 ====
\begin{abstract}
随着城市化进程的加速，交通拥堵已成为深圳等大型城市面临的严峻挑战。传统的交通流量预测方法——如时间序列分析和统计回归——难以捕捉复杂的时空依赖关系，在面对突发状况时预测精度大幅下降。本报告针对深圳市福田区的实际交通场景，设计并实现了一套基于深度学习的智慧交通流量预测系统。系统采用图卷积网络和长短期记忆网络的混合架构，同时建模路网拓扑结构和时序演化规律。实验结果表明，所提方法在预测精度上较基线模型平均提升18.7\%，系统响应延迟低于200ms，满足实时交通管理的需求。

本报告首先对国内外交通流量预测的研究现状进行了系统梳理，然后详细介绍了数据集构建、模型设计与训练优化过程，最后展示了系统部署方案和实际运行效果。报告还讨论了模型的可解释性、边缘部署可行性以及未来扩展方向。
\end{abstract}

% ==== 目录（含插图/表格目录） ====
\tableofcontents

% ==== 引言 ====
\introduction

随着城市化进程的加速，交通拥堵已成为深圳等大型城市面临的严峻挑战。传统的交通流量预测方法难以捕捉复杂的时空依赖关系，在面对突发状况时预测精度大幅下降。本报告针对深圳市福田区的实际交通场景，设计并实现了一套基于深度学习的智慧交通流量预测系统。

引言部分概述了本报告的研究动机、核心问题和整体框架，为后续各章的详细论述奠定基础。报告组织结构如下：第一章介绍研究背景与意义；第二章说明数据来源与预处理方法；第三章详述GCN-LSTM混合模型设计；第四章展示实验结果与性能评估；第五章介绍系统架构与部署方案；第六章总结全文并展望未来方向。

% ============================================================
\chapter{绪论}

\section{研究背景与意义}

交通拥堵已成为制约城市可持续发展的关键瓶颈之一。据深圳市交通运输局统计，2025年福田区高峰期平均车速仅为18.3 km/h，较2019年下降了12.6\%。交通拥堵不仅造成了巨大的经济损失，还加剧了空气污染和噪声问题。

传统的交通管理主要依赖固定配时信号控制和人工调度，缺乏对实时交通状况的动态感知和智能决策能力。近年来，随着传感器技术、5G通信和人工智能的快速发展，智慧交通系统的建设条件日益成熟。

\section{国内外研究现状}

\subsection{传统预测方法}

时间序列分析是最早应用于交通流量预测的方法之一。ARIMA模型及其变体通过捕捉历史数据中的线性趋势和季节性规律进行预测。然而，这些方法难以处理交通流中的非线性突变和突发事件。

\subsection{深度学习方法}

近年来，深度学习技术在交通预测领域展现了巨大潜力。循环神经网络（RNN）及其变体——特别是长短期记忆网络——能够建模交通流的长程时序依赖。图卷积网络则被引入来处理路网的非欧几里得结构。

\section{报告组织结构}

本报告共分为六章，组织结构如下：

\begin{enumerate}
    \item \textbf{第一章 绪论}：介绍研究背景、现状和组织结构
    \item \textbf{第二章 数据集与预处理}：说明数据来源、特征工程和数据集构建
    \item \textbf{第三章 模型设计}：详述 GCN-LSTM 混合架构设计
    \item \textbf{第四章 实验与评估}：展示实验结果和性能分析
    \item \textbf{第五章 系统实现}：介绍系统架构和部署方案
    \item \textbf{第六章 总结与展望}：总结工作并提出未来方向
\end{enumerate}

% ============================================================
\chapter{数据集与预处理}

\section{数据来源}

本报告使用的交通数据来自以下渠道：

\begin{itemize}
    \item \textbf{出租车GPS轨迹数据}：深圳市2019年全量出租车GPS轨迹，采样间隔15秒
    \item \textbf{卡口过车数据}：福田区125个主要路口的车辆通行记录
    \item \textbf{路网拓扑数据}：OpenStreetMap导出的路网结构，包含节点和边信息
    \item \textbf{天气数据}：深圳市气象局提供的逐小时气象观测数据
\end{itemize}

\section{特征工程}

\subsection{时空特征构建}

交通流量具有显著的时空特性。我们以5分钟为时间窗口，为每个路口构建以下特征：

\begin{itemize}
    \item 车流量（veh/5min）
    \item 平均车速（km/h）
    \item 车辆密度（veh/km）
    \item 相邻路口流量差值
\end{itemize}

\subsection{外部特征融合}

除交通数据外，我们还将外部因素作为辅助特征输入模型：

\begin{table}[htbp]
\centering
\begin{tabular}{lcc}
\toprule
\textbf{特征类别} & \textbf{维数} & \textbf{说明} \\
\midrule
时间特征 & 7 & 小时、星期、是否节假日等独热编码 \\
天气特征 & 4 & 温度、湿度、降水量、风速 \\
POI特征  & 10 & 周边设施类型的密度分布 \\
\bottomrule
\end{tabular}
\caption{外部特征说明}
\label{tab:features}
\end{table}

% ============================================================
\chapter{模型设计}

\section{整体架构}

本系统采用 GCN-LSTM 混合架构，如图\ref{fig:architecture}所示。GCN负责提取路网的空间拓扑特征，LSTM负责建模时序演化规律。

\begin{figure}[htbp]
    \centering
    \colorbox{gray!10}{\parbox{0.85\textwidth}{\centering\vspace{2cm}
        {\large\color{gray} [GCN-LSTM 混合架构示意图]}\\[0.5cm]
        \small 输入 → GCN（空间编码）→ LSTM（时序建模）→ 全连接层 → 预测输出
    \vspace{2cm}}}
    \caption{GCN-LSTM 交通流量预测模型架构}
    \label{fig:architecture}
\end{figure}

\section{空间建模：图卷积网络}

将路网建模为有向图 $G = (V, E, A)$，其中 $V$ 为路口集合，$E$ 为道路集合，$A \in \mathbb{R}^{N \times N}$ 为邻接矩阵。

图卷积层的计算过程为：

\begin{equation}
H^{(l+1)} = \sigma\left(\tilde{D}^{-\frac{1}{2}} \tilde{A} \tilde{D}^{-\frac{1}{2}} H^{(l)} W^{(l)}\right)
\label{eq:gcn}
\end{equation}

其中 $\tilde{A} = A + I$ 为添加自环的邻接矩阵，$\tilde{D}$ 为对应的度矩阵。

\section{时序建模：长短期记忆网络}

在GCN提取每时刻的空间特征后，LSTM负责捕获时序依赖：

\begin{align}
f_t &= \sigma(W_f \cdot [h_{t-1}, x_t] + b_f) \label{eq:lstm_f} \\
i_t &= \sigma(W_i \cdot [h_{t-1}, x_t] + b_i) \label{eq:lstm_i} \\
o_t &= \sigma(W_o \cdot [h_{t-1}, x_t] + b_o) \label{eq:lstm_o} \\
c_t &= f_t \odot c_{t-1} + i_t \odot \tanh(W_c \cdot [h_{t-1}, x_t] + b_c) \label{eq:lstm_c} \\
h_t &= o_t \odot \tanh(c_t) \label{eq:lstm_h}
\end{align}

\section{训练策略}

\subhead{损失函数}采用均方误差作为主要损失函数，同时引入平滑正则化项防止过拟合：

\begin{equation}
\mathcal{L} = \frac{1}{N} \sum_{i=1}^{N} (y_i - \hat{y}_i)^2 + \lambda \sum_{i=2}^{N} |\hat{y}_i - \hat{y}_{i-1}|
\label{eq:loss}
\end{equation}

\subhead{优化器配置}使用Adam优化器，初始学习率0.001，训练200个epoch，batch size为64。学习率采用余弦衰减策略。

% ============================================================
\chapter{实验与评估}

\section{实验设置}

实验在配备NVIDIA RTX 3090 GPU和128GB RAM的服务器上完成。数据集按时间顺序划分为训练集（70\%）、验证集（15\%）和测试集（15\%）。

\section{对比实验}

\subsection{基准模型}

选取以下经典方法作为对比基线：
\begin{itemize}
    \item \textbf{HA}：历史平均法，取过去12周的同期均值作为预测
    \item \textbf{ARIMA}：自回归积分滑动平均模型
    \item \textbf{SVR}：支持向量回归
    \item \textbf{FC-LSTM}：全连接LSTM网络
    \item \textbf{STGCN}：时空图卷积网络
\end{itemize}

\subsection{实验结果}

表\ref{tab:results}展示了各方法在三个预测指标上的对比结果。

\begin{table}[htbp]
\centering
\begin{tabular}{lccc}
\toprule
\textbf{方法} & \textbf{MAE} & \textbf{RMSE} & \textbf{MAPE(\%)} \\
\midrule
HA      & 18.42 & 28.17 & 21.35 \\
ARIMA   & 14.38 & 22.56 & 16.82 \\
SVR     & 13.21 & 20.94 & 15.47 \\
FC-LSTM & 11.05 & 17.83 & 12.91 \\
STGCN   & 9.87  & 15.62 & 11.34 \\
\midrule
\textbf{本文方法} & \textbf{7.92} & \textbf{12.45} & \textbf{9.18} \\
\bottomrule
\end{tabular}
\caption{福田区交通流量预测结果对比}
\label{tab:results}
\end{table}

从表\ref{tab:results}可以看出，本文提出的GCN-LSTM混合模型在所有指标上均取得了最优结果，MAE较最佳基线模型STGCN降低了19.8\%。

\section{消融实验}

为验证各模块的贡献，进行了消融实验，结果如表\ref{tab:ablation}所示。

\begin{table}[htbp]
\centering
\begin{tabular}{cccc}
\toprule
\textbf{GCN} & \textbf{LSTM} & \textbf{外部特征} & \textbf{MAE} \\
\midrule
              & \checkmark & \checkmark & 12.31 \\
\checkmark    &            & \checkmark & 11.05 \\
\checkmark    & \checkmark &            & 9.14  \\
\checkmark    & \checkmark & \checkmark & \textbf{7.92} \\
\bottomrule
\end{tabular}
\caption{消融实验结果}
\label{tab:ablation}
\end{table}

结果表明，GCN和LSTM各自都有显著贡献，外部特征的引入进一步提升了预测精度。

% ============================================================
\chapter{系统实现}

\section{系统架构}

系统采用\"数据层-模型层-应用层\"三层架构。数据层负责实时数据采集和预处理；模型层运行GCN-LSTM推理引擎；应用层提供Web可视化界面和API接口。

\section{部署方案}

模型部署在高性能计算中心的GPU服务器上，通过Docker容器化部署确保环境一致性。系统对外提供RESTful API接口，支持深圳市交通运输局现有系统的无缝对接。

\section{性能测试}

在并发1000请求/秒的压力测试下，系统P99延迟为187ms，CPU利用率62\%，满足实时交通管理的服务等级协议要求。

% ============================================================
\chapter{总结与展望}

\section{工作总结}

本报告针对深圳福田区交通流量预测的实际需求，设计并实现了基于GCN-LSTM混合架构的智慧交通预测系统。主要贡献包括：

\begin{enumerate}
    \item 构建了融合多源数据的福田区交通数据集
    \item 提出了GCN-LSTM混合模型，在预测精度上较基线方法提升18.7\%
    \item 完成了系统的工程化部署，验证了实时应用的可行性
\end{enumerate}

\section{未来展望}

未来的研究方向包括：(1) 引入Transformer架构处理更长程的时序依赖；(2) 探索联邦学习解决跨区域数据隐私问题；(3) 集成强化学习实现信号灯自适应控制。

% ==== 基金信息 ====
\printfunding

% ==== 参考文献 ====
\begin{thebibliography}{12}

\bibitem{lstm} Hochreiter S, Schmidhuber J. Long short-term memory[J]. Neural Computation, 1997, 9(8): 1735-1780.

\bibitem{gcn} Kipf T N, Welling M. Semi-supervised classification with graph convolutional networks[C]. ICLR, 2017.

\bibitem{stgcn} Yu B, Yin H, Zhu Z. Spatio-temporal graph convolutional networks: A deep learning framework for traffic forecasting[C]. IJCAI, 2018.

\bibitem{dcrnn} Li Y, Yu R, Shahabi C, et al. Diffusion convolutional recurrent neural network: Data-driven traffic forecasting[C]. ICLR, 2018.

\bibitem{astgcn} Guo S, Lin Y, Feng N, et al. Attention based spatial-temporal graph convolutional networks for traffic flow forecasting[C]. AAAI, 2019.

\bibitem{gwn} Wu Z, Pan S, Long G, et al. Graph WaveNet for deep spatial-temporal graph modeling[C]. IJCAI, 2019.

\bibitem{sz_plan} 深圳市交通运输局. 深圳市综合交通"十四五"规划[R]. 深圳, 2021.

\bibitem{traffic_review} Jiang W, Luo J. Graph neural network for traffic forecasting: A survey[J]. Expert Systems with Applications, 2022, 207: 117921.

\bibitem{alibaba} 阿里云计算有限公司. 城市大脑：数据智能驱动的城市治理[R]. 杭州, 2020.

\end{thebibliography}

\end{document}
